\documentclass[11pt]{article}
\usepackage[a4paper,margin=25mm]{geometry}
\usepackage[T1]{fontenc}
\usepackage[utf8]{inputenc}
\usepackage{lmodern,microtype,amsmath,amssymb,amsthm,booktabs,graphicx,natbib}
\usepackage[colorlinks=true,allcolors=blue,breaklinks=true]{hyperref}
\usepackage{url}
\newtheorem{proposition}{Proposition}
\newcommand{\E}{\mathbb{E}}
\newcommand{\R}{\mathbb{R}}
\newcommand{\tr}{\operatorname{tr}}
\newcommand{\norm}[1]{\left\lVert #1\right\rVert}
\title{When does force-label denoising change an invariant potential?\\[3pt]
\large Symmetry, covariance and paired-DFT evidence}
\author{Sinan G\"okmen\\Department of Chemistry, Istanbul Technical University\\
\texttt{gokmens23@itu.edu.tr}}
\date{11 September 2026}
\begin{document}
\maketitle
\begin{abstract}
Removing numerical error from force labels is a plausible way to improve interatomic
potentials, but a decrease in label error need not change learning. For an invariant
energy model trained with squared force loss, we show that orthogonal removal of
net force and torque changes the objective only by a parameter-independent constant.
Useful correction must reach the complementary, symmetry-allowed subspace. A
conditional-noise formulation identifies the cross-covariance that enables this
correction and gives an exact excess-risk identity for its estimation. We examine
these ideas on 2,000 publicly available paired DFT configurations using molecular-
formula-grouped evaluation. On the AIMNet2 sample, element-weighted constraint
correction reduces force RMSE from 4.7979 to 4.4800 meV/\AA; a spectral bulk-pooling
heuristic reaches 4.4909 meV/\AA. The latter does not establish superiority over ridge
and is inferior to the element-weighted baseline. On Transition1x, the tested methods
make little difference to aggregate error. A separate Gaussian experiment illustrates
when spiked-covariance shrinkage can help, with an explicitly known noise floor.
These results identify a necessary distinction between symmetry repair, useful label
correction and improved model prediction; the last is not established here.
\end{abstract}

\section{The question behind force denoising}
Force labels influence both the fitting and assessment of machine-learned interatomic
potentials. Paired recalculations by \citet{kuryla2025} provide an opportunity to study
numerical force discrepancies directly. However, a model cannot learn every component
of a force vector. In particular, derivatives of a translation- and rotation-invariant
energy obey zero-net-force and zero-net-torque conditions. A preprocessing method can
therefore reduce a visible label discrepancy while leaving every training gradient
unchanged.

This observation motivates three distinct questions: which force errors are invisible
to the chosen learner, whether observable constraint violations predict the remaining
error, and how covariance-estimation error affects such prediction. The calculations
below use elementary orthogonal projection and linear prediction; neither mathematical
tool is claimed as new. The contribution investigated here is their joint use as an
evaluation criterion for force-label denoising, accompanied by a reproducible paired-data
comparison. Priority for this formulation has not been established by an exhaustive
literature review.

Learning conservative forces directly from an energy-related function space is an
established strategy \citep{chmiela2017}. Coordinate-denoising pretraining
\citep{wang2023} and weak supervision using energy--force relations
\citep{takamoto2024} address related learning problems. The present intervention acts
on numerical force labels; it does not perturb molecular coordinates or infer a
binding pose. Random matrix theory (RMT) is relevant to estimating high-dimensional
covariance spectra \citep{marchenko1967,donoho2017}, while covariance shrinkage also
admits strong nonspectral baselines \citep{chen2009}. Its utility must be assessed on
the resulting correction, rather than inferred from a smaller covariance error alone.

\section{An exact limitation of symmetry projection}
Consider an isolated molecule with positions $R=(r_1,\ldots,r_N)$ and a differentiable
energy $E_\theta(R)$ invariant under translations and rotations. Stack its forces in
$f_\theta=-\nabla_R E_\theta\in\R^{3N}$. Define the constraint operator
\begin{equation}
 A_R f=\begin{pmatrix}\sum_i f_i\\\sum_i (r_i-\bar r)\times f_i\end{pmatrix},
 \qquad \bar r=N^{-1}\sum_i r_i.
\end{equation}
Differentiating the two invariances yields $A_R f_\theta=0$. Let $U_R$ have orthonormal
columns spanning $\operatorname{range}(A_R^T)$, and write
\begin{equation}
 Q_R=U_RU_R^T,\qquad P_R=I-Q_R.
\end{equation}
These definitions also apply to linear or degenerate geometries if the actual rank
is used. The word ``tangent'' below abbreviates the symmetry-allowed subspace
$\ker A_R$; it does not assert that every vector field taking values there is
globally conservative.

\begin{proposition}[Projection does not change squared-loss learning]
For any fixed force label $y$,
\begin{equation}\label{eq:noop}
 \norm{f_\theta-y}^2=\norm{f_\theta-P_Ry}^2+\norm{Q_Ry}^2.
\end{equation}
Consequently, replacing $y$ by $P_Ry$ leaves parameter gradients, Hessians and the
set of minimizers unchanged for a squared force loss with fixed scalar weights per
configuration. The conclusion remains valid if the same energy loss and parameter
regularizer are added to both objectives.
\end{proposition}
\begin{proof}
The vector $f_\theta-P_Ry$ belongs to $\ker A_R$, and $Q_Ry$ is orthogonal to that
subspace. Their squared norms add. The last term in \eqref{eq:noop} is independent
of $\theta$.
\end{proof}

This is a statement about the objective, not every possible training program.
Absolute-loss stopping rules, label-dependent sampling, robust component losses,
learned loss weights and unequal atomwise metrics may react differently. The
statement also requires the model itself to obey the invariances. Equivariance of
a direct force predictor alone does not guarantee zero net force or torque.

An additional restriction appears for any linear conservative model $f_\theta=J\theta$.
For ridge fitting with $H=J^TJ+\lambda I$, replacing targets $y$ by $y-d$ changes
the fitted coefficients by
\begin{equation}\label{eq:response}
 \widehat\theta_{y-d}-\widehat\theta_y=-H^{-1}J^Td.
\end{equation}
Thus $P_Rd\ne0$ is necessary to escape the symmetry limitation but is not sufficient
to affect a particular learner: a correction may still be orthogonal to its design
space. Nor does a nonzero change in parameters guarantee lower test error.

\section{A covariance condition for useful correction}
For this section fix a geometry and an orthonormal complementary basis $V_R$, with
$V_RV_R^T=P_R$. Suppose an ideal reference force $f_*$ satisfies $A_Rf_*=0$ and
$y=f_*+\epsilon$. Assume conditional zero-mean noise with covariance $\Sigma$;
all expectations may be interpreted conditional on the geometry and acquisition
settings. Define
\begin{equation}
 \epsilon_t=V_R^T\epsilon,\quad \epsilon_n=U_R^T\epsilon,
 \qquad
 \begin{pmatrix}\Sigma_{tt}&\Sigma_{tn}\\\Sigma_{nt}&\Sigma_{nn}\end{pmatrix}
 =\begin{pmatrix}V_R^T\\U_R^T\end{pmatrix}\Sigma(V_R,U_R),
\end{equation}
and assume $\Sigma_{nn}\succ0$. Constraint violations reveal
$z=U_R^Ty=\epsilon_n$. A linear corrected label is
\begin{equation}\label{eq:correct}
 \widetilde f_B=P_Ry-V_RBz.
\end{equation}
This uses a measured error indicator to correct an error component that otherwise
survives symmetry projection.

\begin{proposition}[Predictable error and estimation penalty]
Among deterministic linear maps $B$, the minimum expected squared error in
\eqref{eq:correct} is attained at $B_* =\Sigma_{tn}\Sigma_{nn}^{-1}$. Writing
$\mathcal R(B)=\E\norm{\epsilon_t-B\epsilon_n}^2$, one has
\begin{align}
 \mathcal R(B_*)&=\tr\bigl(\Sigma_{tt}-\Sigma_{tn}\Sigma_{nn}^{-1}\Sigma_{nt}\bigr),\label{eq:oracle}\\
 \mathcal R(B)-\mathcal R(B_*)&=
 \norm{(B-B_*)\Sigma_{nn}^{1/2}}_F^2.\label{eq:risk}
\end{align}
The population gain over orthogonal projection is
\begin{equation}\label{eq:gain}
 \mathcal R(0)-\mathcal R(B_*)=
 \norm{\Sigma_{tn}\Sigma_{nn}^{-1/2}}_F^2.
\end{equation}
\end{proposition}
\begin{proof}
Expand the quadratic expectation as
$\tr\Sigma_{tt}-2\tr(B\Sigma_{nt})+\tr(B\Sigma_{nn}B^T)$ and complete the square.
Positive definiteness gives uniqueness. Substitution of $B=0$ gives \eqref{eq:gain}.
\end{proof}

For a random estimator $\widehat B$ learned from independent calibration data,
\eqref{eq:risk} holds conditional on that estimate, and its expectation adds the
calibration penalty. If joint Gaussianity holds, $B_*z$ is also the conditional mean
of $\epsilon_t$ given $z$, hence optimal among all square-integrable predictors of
that indicator. Without Gaussianity only the linear optimality claim is made.
With systematic conditional bias, the corresponding conditional means must also
be estimated; numerical DFT discrepancies need not be zero-mean Gaussian noise.

Two consequences are immediate. If $\Sigma_{tn}=0$, no linear correction based
only on constraint violations improves population risk. If the cross-covariance is
nonzero, its predictable gain must exceed the finite-calibration penalty. A large
force-error variance or a large spectral spike by itself does not imply that gain.

For a plug-in estimator $\widehat B=\widehat\Sigma_{tn}\widehat\Sigma_{nn}^{-1}$,
suppose both normal-block minimum eigenvalues exceed $a>0$. Then
\begin{equation}\label{eq:bound}
 \norm{\widehat B-B_*}_F\le
 a^{-1}\norm{\widehat\Sigma_{tn}-\Sigma_{tn}}_F+
 a^{-2}\norm{\Sigma_{tn}}_2\norm{\widehat\Sigma_{nn}-\Sigma_{nn}}_F.
\end{equation}
This follows by adding and subtracting
$\Sigma_{tn}\widehat\Sigma_{nn}^{-1}$ and using the inverse-difference identity.
Together with \eqref{eq:risk}, it exposes both cross-block error and instability of
the normal-block inverse. Shrinking unrelated tangent variances need not help.

In Cartesian coordinates the population linear correction can equivalently be
written
\begin{equation}\label{eq:oblique}
 \widetilde f=y-\Sigma U_R(U_R^T\Sigma U_R)^{-1}U_R^Ty.
\end{equation}
It satisfies the constraints. For $\Sigma=\sigma^2I$, it reduces to $P_Ry$ and has
the learning limitation of Proposition 1. The map in \eqref{eq:oblique} is a
classical covariance-weighted constraint correction, not a newly invented estimator.

\section{A paired-DFT pilot}
\subsection{Data and separation of validation from test}
We use the public paired-force files accompanying \citet{kuryla2025}, pinned to
commit
\begin{center}\small\texttt{322a41229f7707a58626a7c03ce850d67920fad0}.\end{center}
Their tightly recomputed
ORCA forces are operational references, not exact electronic-structure truth.
The AIMNet2 and Transition1x files each contain 1,000 configurations; direct use
of their force columns reproduces the reported force RMSDs. These calculations
compare numerical setups at a fixed nominal level of theory; they do not establish
removal of exchange-correlation functional bias.

An audit of two additional upstream files did not reproduce the stated comparisons
under the attempted unit and dispersion conventions. ANI-1x and SPICE are therefore
excluded from fitted experiments. Appendix A records the unresolved checks without
attributing their cause to the source publication.

Entire molecular formulas are assigned deterministically to five outer folds using
a SHA-256 hash. For each fold, one group is test, the next is validation, and the
remaining three are calibration. This produces one out-of-fold prediction per frame.
There are 977 formula groups in AIMNet2 and 122 in Transition1x. Calibration sets
contain 574--627 and 558--655 configurations, respectively. Shared fixed-composition
conformers cannot cross partitions, but this is not a scaffold split or a guaranteed
reaction-family split. All configurations, including large-error outliers, are kept.

\subsection{Corrections and baselines}
The untrained baselines are raw forces, removal of the mean force, and orthogonal
force/torque projection. An element-weighted baseline uses \eqref{eq:oblique} with
a diagonal covariance whose per-atom Cartesian variance is estimated from that
element's calibration discrepancies. Unseen elements use the pooled calibration
variance. A small floor prevents zero variance.

The learned corrections use an equivariant design $X_R(y)\in\R^{3N\times48}$.
Its columns are the projected products of an element indicator, an invariant
environment scalar and a vector marker. The eight element channels are H, C, N,
O, F, S, Cl and other. The scalars are $1$, $s_i/(1+s_i)$ and $c_i/(1+c_i)$, where
$s_i=\norm{r_i-\bar r}$ in \AA\ and
$c_i=\sum_{j\ne i}\exp(-\norm{r_i-r_j}^2/(4\,\mathring{\mathrm A}^{2}))$.
The markers are the translational component of $y$ and the remaining rotational
component of $Q_Ry$. Each resulting column is projected by $P_R$.
The prediction is $P_Ry-X_R(y)\beta$, where $\beta$ is fitted to
$P_R(y-f_{\mathrm{ref}})$ using calibration frames only. These corrected labels
obey the constraints but need not form a conservative field across geometries.

We compare ridge regression, spectral truncation and an MP-inspired bulk-pooling
heuristic on the diagonally normalized feature Gram matrix. Each method selects
its regularization using validation MSE and includes zero correction as a candidate.
For the heuristic the nominal edge is $(1+\sqrt{p/m})^2$, with $p$ active columns
and $m$ calibration configurations. Eigenvalues below the edge are partly pooled
to their mean; blending and a ridge floor are selected on validation data.
Atomic components are dependent and the features are structured, so this Gram
matrix is not demonstrated to satisfy a Wishart or isotropic MP model. The edge
is explicitly a heuristic here, and is different from the justified Gaussian setup
in Section 5. No asymptotic optimality is asserted for the molecular correction.

\subsection{Results}
\begin{table}[t]
\centering\small
\begin{tabular}{lrrrr}
\toprule
&\multicolumn{2}{c}{Pooled force RMSE}&\multicolumn{2}{c}{Median frame RMSE}\\
Method&AIMNet2&Transition1x&AIMNet2&Transition1x\\
\midrule
Raw & 4.7979 & 26.9007 & 4.3459 & 3.8608 \\
Translation & 4.7143 & 26.8942 & 4.2986 & 3.8150 \\
Force + torque & 4.6092 & 26.8853 & 4.2338 & 3.7509 \\
Element weighted & 4.4800 & 26.8872 & 4.1333 & 3.8059 \\
Ridge correction & 4.4940 & 26.8836 & 4.1287 & 3.7114 \\
Spectral truncation & 4.4951 & 26.8855 & 4.1314 & 3.7192 \\
MP bulk heuristic & 4.4909 & 26.8833 & 4.1289 & 3.7239 \\
\bottomrule
\end{tabular}

\caption{Out-of-fold force-label discrepancies in meV/\AA. The pooled metric weights
Cartesian components equally. No trained potential prediction is included in this table.}
\label{tab:real}
\end{table}

On AIMNet2, 92.28\% of the squared discrepancy lies in the symmetry-allowed
subspace; on Transition1x, the fraction is 99.89\%. The tight references have small
but nonzero forbidden components: RMS values are 0.0220 and 0.0517 meV/\AA,
respectively. The ideal-reference premise of Section 3 is therefore approximate
for the actual paired data.

Table~\ref{tab:real} separates label improvement from downstream learning.
Element weighting lowers AIMNet2 RMSE by 6.63\% relative to raw labels, and by
2.80\% relative to orthogonal projection. The MP heuristic improves on orthogonal
projection by 0.1183 meV/\AA. A paired formula-cluster bootstrap with 2,000 draws
gives a 95\% interval of $[-0.1329,-0.1033]$ meV/\AA\ for that difference.
However, its difference from ridge is $-0.0031$ with interval
$[-0.0069,0.0002]$, and its difference from element weighting is $+0.0109$ with
interval $[0.0046,0.0172]$. Thus the experiment does not establish a special RMT
advantage; the simpler element-weighted method has the lower aggregate error.

\begin{figure}[t]
\centering\includegraphics[width=\linewidth]{figures/real_dft_comparison.pdf}
\caption{Force-label RMSE for the seven methods. Axes begin at zero. Improvements
on AIMNet2 are modest; aggregate Transition1x error changes little.}
\end{figure}

On Transition1x, the MP heuristic changes pooled RMSE by only $-0.0020$ meV/\AA\
relative to orthogonal projection, with interval $[-0.0311,0.0012]$. The large gap
between median-frame and pooled error is consistent with strong tail influence.
Those observations are retained in the primary endpoint. The bootstrap describes
resampling of the existing out-of-fold predictions conditional on the fitted
models; it does not refit every model or include all uncertainty from choosing
these dataset samples and feature families.

\section{Controlled covariance mechanism}
To illustrate the exact risk decomposition separately from the molecular comparison,
we specify a Gaussian calibration experiment with $p=24$ components and six
normal components. This illustration was designed after the real-data comparison;
it is not a preregistered replication or independent evidence of molecular gains.
The three covariances are $I$, $I+10ww^T$ with
$w=(e_1+e_7)/\sqrt2$, and $I+10e_7e_7^T$. They represent no correlation, a
cross-subspace spike and a tangent-only spike, respectively.

For $m\in\{16,32,64,128\}$, we draw 200 independent calibration sets from each
covariance. We compare orthogonal projection, the empirical covariance correction,
OAS \citep{chen2009}, a spiked-model spectral shrinker \citep{donoho2017}, and the
population oracle. Expected test risk is evaluated analytically from the known
covariance using \eqref{eq:risk}; no sampled test noise is needed. The spectral
method is supplied the true isotropic noise floor of one. This is a substantial
oracle advantage and is not available in the molecular experiment.

\begin{figure}[t]
\centering\includegraphics[width=\linewidth]{figures/synthetic_risk.pdf}
\caption{Expected tangent error in three Gaussian models. Curves average 200
independent calibrations; bands are 1.96 standard errors across calibrations.
The spiked spectral estimator receives the known noise floor. A tangent-only
spike carries no normal-to-tangent prediction opportunity despite its large eigenvalue.}
\label{fig:synthetic}
\end{figure}

At $m=32$, the cross-subspace example has per-tangent-component expected MSE
1.2778 for projection, 1.2868 for empirical covariance, 1.1346 for OAS,
1.0761 for the known-floor spectral estimator and 1.0463 for the population oracle.
In contrast, for the tangent-only spike, orthogonal projection and the oracle
both achieve 1.5556; the spectral estimator yields 1.6122. The latter estimator
cannot create information absent from the cross-covariance. These examples show
both a setting where spectral regularization is useful and one where denoising
the covariance does not improve constraint-based force correction.

\section{Implications and limits}
The combination of Propositions 1 and 2 supplies a concrete criterion for a denoising
study: report how much error is removed from the component visible to the learner,
and compare the calibration penalty with the predictable gain. A decrease in net
force is insufficient evidence. A decrease in covariance Frobenius error is also
insufficient unless it improves the relevant conditional map or final prediction.

The molecular pilot supports modest label correction on one dataset, not a more
accurate trained potential. It uses hundreds of tightly recomputed calibration
and validation labels. Any claim of data or compute efficiency must charge those
labels to the method and compare against spending the same budget directly on
potential training or label replacement. An appropriate subsequent test would
hold model architecture and total reference budget fixed, tune all methods on
independent validation data, and evaluate force and energy prediction on tight
references from unseen molecules. Multiple training seeds and learning curves
would be needed before claiming a smaller model or fewer labels suffice.

These data contain molecular configurations, not protein--ligand binding free
energies. The correction does not establish improved docking, affinity ranking,
kinetics or drug discovery. Extending it to such endpoints would require
corresponding experiments and treatment of sampling, solvent and systematic
electronic-structure error. Coordinate noise, force-label error and genuine
thermal variation should not be merged into a single noise model.

\section{Conclusion}
Orthogonal force/torque repair is exactly invisible to squared-loss training of an
invariant energy model. Useful constraint-informed correction requires predictable
error inside the symmetry-allowed subspace, and its success depends on estimating
the relevant cross-covariance. Our paired-data experiment finds modest gains on
AIMNet2 but no advantage of the tested RMT-inspired method over a simpler
element-weighted baseline. The controlled Gaussian calculations delimit a mechanism
for spectral improvement without establishing its dominance in computational
chemistry. Downstream predictive and efficiency claims remain open.

\section*{Data and code}
The original paired data are available at
\url{https://github.com/water-ice-group/datasets_dft_accuracy}.
The accompanying research archive contains the analysis scripts, result tables,
the recorded experiment protocol, immutable download URLs and SHA-256/Git-blob
checksums. Source data are retrieved with the downloader rather than repackaged.
No new electronic-structure calculations were performed. A public release of
this analysis code has not yet been made.

\section*{Generative AI assistance}
ChatGPT/Codex was used extensively in developing the mathematical arguments,
writing and checking the analysis code, running the numerical experiments,
searching the literature, and preparing the text and figures. Numerical results
were generated by the supplied scripts. This disclosure does not represent an
independent verification of the scientific claims.

\appendix
\section{Audit and implementation details}
\paragraph{Unresolved upstream comparisons.}
Using the upstream ANI conversion of Hartree/\AA\ to eV/\AA\ on the available
\texttt{ani\_1x\_tz\_configs.xyz} gives 173.86 meV/\AA\ and near-zero original net
forces; it does not reproduce the reported 33.2 meV/\AA\ comparison. The available
SPICE file contains 998 configurations, a gradient-sign convention and a separate
D3 array. The attempted conversion $-\mathrm{REF\_forces}-\mathrm{dftd3\_forces}$
gives 16.68 meV/\AA, not the reported comparison. These outputs are diagnostic only.
Neither file is accepted into the fitted benchmark. Their precise correspondence
to published computational settings remains unresolved.

\paragraph{Projection.}
Three translation columns and three infinitesimal-rotation columns are constructed
at centered coordinates. Translation columns are divided by $\sqrt N$ and rotation
columns by the centered-coordinate Frobenius norm. An SVD retains singular values
above $10^{-10}$ times the largest. This changes numerical scaling, not the subspace.

\paragraph{Regression.}
For calibration frames, $G=\sum X^TX$ and $b=\sum X^TP_R(y-f_{\mathrm{ref}})$.
Columns with diagonal entries below $10^{-12}\max\operatorname{diag}G$ or the
absolute floor $10^{-28}$ are removed. The remaining Gram matrix is scaled by
$D=\operatorname{diag}(\sqrt{G_{jj}})$, so $H=D^{-1}GD^{-1}$ has unit diagonal.
Ridge penalties are $10^{-6},10^{-5},\ldots,10^2$. Truncation cutoffs relative to
the largest eigenvalue are $10^{-6},10^{-4},10^{-3},10^{-2},0.03,0.1,0.3$.
MP bulk blend strengths are $0.25,0.5,0.75,1$ and ridge floors are
$10^{-6},10^{-3},0.1,1$. No outer test label enters these choices. The
element-variance floor is $10^{-3}$ times the pooled variance, or $10^{-20}$,
whichever is larger.

\paragraph{Known-floor shrinker.}
In the Gaussian illustration, let $\gamma=p/m$ and $\lambda$ be a sample
eigenvalue. Eigenvalues below $(1+\sqrt\gamma)^2$ are replaced by one. Above that
edge, set
\begin{align}
 \ell&=\tfrac12\bigl[\lambda+1-\gamma+
 \sqrt{(\lambda+1-\gamma)^2-4\lambda}\bigr],\\
 c^2&=\frac{1-\gamma/(\ell-1)^2}{1+\gamma/(\ell-1)},\qquad
 \eta(\lambda)=1+(\ell-1)c^2.
\end{align}
This applies the spiked-model Frobenius shrinkage formula
\citep{donoho2017}; no finite-sample optimality for conditional force risk is claimed.

\paragraph{Verification.}
Analytic conservative pair-potential columns were used to check invariance and
\eqref{eq:noop} on 80 configurations. Maximum normal-equation disagreement was
$5.33\times10^{-15}$, and maximum relative squared-loss identity error was
$4.67\times10^{-16}$. Rotation/translation and atom-permutation checks of the
correction design were below $1.55\times10^{-15}$ in absolute value. Across
12,000 synthetic method results the risk identity residual was below
$7.11\times10^{-15}$. These numerical checks support implementation consistency;
the propositions themselves are proved algebraically above.

\bibliographystyle{plainnat}
\bibliography{references}
\end{document}
